% ================================================================
% 《分析讲义》中译本 —— 类文件之上的兼容层（mathtranslation.cls v3.1）
% ----------------------------------------------------------------
% mathtranslation.cls v3.1 基于 ctexbook，已原生提供 \chapter/\part/\appendix、
% 元信息命令 \BookTitleCN 等、\newterm/\addterm 术语系统、\printterminology、
% \headfont 标题字体、\makecover 封面、\longproof、\exercises/\answers、
% mtfang/mtsong 定理样式与 envenddiamond 结束标记。
%
% 因此本文件只保留 v3.1 未提供、且本书确需的部分；不改动正文内容：
%   1. 定理类环境共享同一计数器，编号“节号.序号”（与原书一致）。
%   2. \mtsubsec{标题} / \mtsubsecn{数字}{标题}：节内带编号/独立编号子标题。
%   3. 未编号环境 theorem* / definition* / ...（原书未编号对象用）。
%   4. \thmref 等按类型引用命令。
%   5. mybibliography 兼容环境（v3.0 起类文件只提供 \makebibliography）。
% ================================================================
% ---------- 1. 元信息命令 ----------
% v3.1 已把 \BookTitleCN / \OriginalAuthor / \Translator 等全部暴露为公开命令
% （见类文件第 7 节），旧版为 v1.2 写的 \MT@ 别名层在此整体删除。

% ---------- 2. 共享定理计数器 ----------
% v3.1 把定理类环境各自注册为 [section] 计数，各环境互相独立。本书要求它们
% 共用同一计数器，编号显示为“节号.序号”，且序号全书连续（与原书一致）。
% 故在此解除原定义并重建。
%
% 只 \let 环境命令，不释放 \c@theorem 等计数器本身：v3.1 已把这些计数器登记
% 进 \cl@section 的重置链，销毁计数器会让后续 \section 报
% “No counter 'definition' defined.”（报错位置会被误指到 \section 一行）。
\makeatletter
\let\theorem\relax\let\endtheorem\relax
\let\definition\relax\let\enddefinition\relax
\let\lemma\relax\let\endlemma\relax
\let\corollary\relax\let\endcorollary\relax
\let\proposition\relax\let\endproposition\relax
\let\remark\relax\let\endremark\relax
\let\example\relax\let\endexample\relax
\let\problem\relax\let\endproblem\relax
\makeatother

% 共享计数器：全书连续（不随 chapter/section 重置），显示“节号.序号”。
\newcounter{mtthm}
\renewcommand{\themtthm}{\arabic{section}.\arabic{mtthm}}

% 定理样式直接沿用类文件的 mtfang / mtsong（例、问题用宋体），不再另行定义。
\theoremstyle{mtfang}
\newtheorem{theorem}[mtthm]{定理}
\newtheorem{definition}[mtthm]{定义}
\newtheorem{lemma}[mtthm]{引理}
\newtheorem{corollary}[mtthm]{推论}
\newtheorem{proposition}[mtthm]{命题}
\newtheorem{remark}[mtthm]{注}
\theoremstyle{mtsong}
\newtheorem{example}[mtthm]{例}
\newtheorem{problem}[mtthm]{问题}

% 未编号版本（原书未编号的定义、定理、注等）。声明顺序照旧，因此它们沿用
% 上面生效的 mtsong 样式（宋体），与原书观感一致。
\newtheorem*{theorem*}{定理}
\newtheorem*{definition*}{定义}
\newtheorem*{lemma*}{引理}
\newtheorem*{corollary*}{推论}
\newtheorem*{proposition*}{命题}
\newtheorem*{remark*}{注}
\newtheorem*{example*}{例}
\newtheorem*{problem*}{问题}
\newtheorem*{conjecture*}{猜想}

% 类文件已用 \AtEndEnvironment 给 example/remark 挂了菱形结束标记（钩子存于
% \@end@<env>@hook，由通用 \end 调用，故不必也不可重复挂载，否则会出现双菱形）。
% 下面只为新定义的带 * 未编号环境补挂，保持与编号版本外观一致。
\AtEndEnvironment{example*}{\envenddiamond}
\AtEndEnvironment{remark*}{\envenddiamond}

% 节标题显示为“节号.”（原书样式，不带章号）。
\renewcommand{\thesection}{\arabic{section}.}

% ---------- 3. 节内子标题 ----------
% 带编号子标题: 推进共享计数器并显示 "节号.序号 标题".
\newcommand{\mtsubsec}[1]{%
  \par\medskip\noindent\refstepcounter{mtthm}%
  {\headfont\arabic{section}.\arabic{mtthm}\quad #1}\par\smallskip
}
% 独立编号子标题: 只显示 "数字. 标题", 不推进计数器.
% 用法: \mtsubsecn{数字}{标题}(数字为 0 或空时只显示标题).
\newcommand{\mtsubsecn}[2]{%
  \par\medskip\noindent%
  {\headfont\ifx\relax#1\relax\else #1.\quad\fi#2}\par\smallskip
}

% ---------- 4. 按类型区分的引用命令 ----------
% 所有定理类环境共享同一计数器, 故 \autoref 无法区分类型(一律显示"定理").
% 引用时用下列命令: \thmref{...} → "定理 N", \lemref{...} → "引理 N" 等.
\newcommand{\thmref}[1]{定理~\ref{#1}}
\newcommand{\lemref}[1]{引理~\ref{#1}}
\newcommand{\defref}[1]{定义~\ref{#1}}
\newcommand{\propref}[1]{命题~\ref{#1}}
\newcommand{\corref}[1]{推论~\ref{#1}}
\newcommand{\remref}[1]{注~\ref{#1}}
\newcommand{\exref}[1]{例~\ref{#1}}
\newcommand{\probref}[1]{问题~\ref{#1}}

% ---------- 5. 术语系统 ----------
% v3.1 的 \newterm{键}{中文}{英文} 与旧版行为完全一致（首次出现自动注册、
% 自动取页码、自动去重，并输出到 \MTTermList），故删除此处为 v1.2 写的
% 重新实现，直接使用类文件版本。

% ---------- 6. 公式/图表编号 ----------
% ctexbook（book 基类）默认把 equation / figure / table 按章编号（\theequation =
% 章号.序号）。原书与旧版 ctexart 类均为全书连续编号（公式 (1)--(446)，见"译者
% 说明"），故解除按章重置、恢复连续编号，保持与原文及既有成果一致。
\makeatletter
\@removefromreset{equation}{chapter}
\renewcommand{\theequation}{\arabic{equation}}
\@removefromreset{figure}{chapter}
\renewcommand{\thefigure}{\arabic{figure}}
\@removefromreset{table}{chapter}
\renewcommand{\thetable}{\arabic{table}}
\makeatother

% ---------- 6b. mybibliography 兼容环境 ----------
% v1.0 提供 mybibliography（参考文献列表，标题由正文手动 \section* 给出）。
% v1.2 已删除该环境；此处按 v1.2 的 thebibliography 的列表部分重建（不含标题），
% 供 partapx.tex 的"参考文献"使用。
\makeatletter
\newenvironment{mybibliography}[1]{%
  \begingroup
  \list{\@biblabel{\@arabic\c@enumiv}}%
       {\settowidth\labelwidth{\@biblabel{#1}}%
        \leftmargin\labelwidth
        \advance\leftmargin\labelsep
        \@openbib@code
        \usecounter{enumiv}%
        \let\p@enumiv\@empty
        \renewcommand\theenumiv{\@arabic\c@enumiv}}%
  \sloppy
  \clubpenalty4000
  \@clubpenalty \clubpenalty
  \widowpenalty4000%
  \sfcode`\.\@m
}{%
  \def\@noitemerr{\@latex@warning{Empty bibliography environment}}%
  \endlist
  \endgroup
}
\makeatother

% ---------- 7. 封面 ----------
% v3.1 的 \makecover 已按制作规范印出“译者 xx 翻译及重排”，并逐项拼接
% 版次·出版社·年份（空字段自动跳过），旧版为 v1.2 重写的封面宏在此整体删除。
